% OPTION 1 - draft subsection for insertion after "Robustness & Stability" (before
% "Alternate System"). Budget: ~2/3 IEEE column including Fig. \ref{fig:YH-AW-18}.
%
% Figures and numbers: ../make_YH-AW-18.py and ../make_YH-vs-H-cond.py (or the MATLAB
% script ../antiwindup_variants.m, sections A and B) from Eqs. (\ref{eq:Gp_numerical}),
% (\ref{eq:Gch}) and (\ref{eq:Gcyh}). The torque limit is illustrative (see text).
%
% Framing note (2026-09-16): the conditioned comparison showed that with full-state
% conditioning the H-inf and Youla-H controllers saturate and recover IDENTICALLY -
% Youla-H preserves the transmission zeros, and those zeros are the saturated-mode
% spectrum. The text below therefore claims a well-posed realization and a preserved
% DC benefit, NOT better anti-windup. Figure YH-vs-H-cond-18 is optional; it is
% referenced in the last paragraph and can be dropped with that paragraph if the
% budget is tight.

\subsection{Discretization and Anti-Windup}
\label{sec:antiwindup}
Because the integrator Youla-H plants in $G_{CYH}$ is an \emph{exact} pole at $s=0$,
the controller separates uniquely into an integrator and a strictly stable remainder,
\begin{equation} \label{eq:split}
G_{CYH}(s) = \frac{k_I}{s} + R(s), \qquad
k_I = K_{YH}\frac{\prod_{i} z_i}{\prod_{p_j\neq 0} p_j}.
\end{equation}
No such split exists for $G_{CH}$, whose integral action is buried in the slow pole
pair of Equation \ref{eq:Gch}. The split makes the digital realization well posed:
the integrator is discretized by the trapezoidal rule as an explicit scalar state,
and $R(s)$ is balance-truncated (which requires a stable system) and discretized by
Tustin's method \citep{tustin1947method}. The explicit integrator state also admits the
standard PI
anti-windup schemes, back-calculation or conditional integration, directly.

Figure \ref{fig:YH-AW-18} applies this to the drivetrain system with a $5$\,m/s step
and an illustrative limit $|T_e|\le 0.1$\,N$\cdot$m, chosen to rail the actuator for
about $2$\,s. With a clamp alone, $G_{CH}$ and $G_{CYH}$ alike wind up, overshooting
$71\%$ with a $10$\,s relay-like recovery. Back-calculating into the integrator alone
releases the rail but leaves an $18\%$ standing error: here $k_I = 1.8\times10^{-3}$
while $R(0) = 42$, so the lag states of $R$ hold most of the low-frequency gain and
wind up on their own. Conditioning the full state with the general Hanus law
\citep{hanus1987conditioning}, $x_{k+1} = A_d x_k + B_d e_k + L\,(u_{\mathrm{sat},k}-u_k)$,
leaves the linear controller unchanged off the clamp and de-winds every state on it:
$1.3\%$ overshoot, settled $2.2$\,s after the step. $L$ is pole-placed rather than
the self-conditioned $B_d/D_d$, because Tustin's method gives Equation \ref{eq:split}
a $(z+1)$ numerator factor that puts a saturated-mode eigenvalue at exactly $z=-1$.

Full-state conditioning applies equally to $G_{CH}$, and the two controllers then
saturate and recover identically: the self-conditioned saturated-mode eigenvalues
are the controller's transmission zeros, which the Youla-H gain adjustment leaves
unchanged. Youla-H therefore offers no anti-windup advantage of its own. What it
offers is that its benefit survives discretization and saturation intact. On the
conditioned step both controllers settle in $2.2$\,s, but $G_{CH}$ retains a
residual error of $-5.7\times10^{-4}$\,m/s where $G_{CYH}$ retains none, and on a
slow unsaturated ramp the $G_{CH}$ error grows to $-0.6$\,mm/s over $200$\,s while
$G_{CYH}$ tracks exactly (Figure \ref{fig:YH-vs-H-cond-18}).

\begin{figure}[H]
    \includegraphics[width=1\linewidth]{antiwindup_study/figures/YH-AW-18.pdf}
    \caption{Saturated step response of the drivetrain system (reference in light
    gray): clamp only, integrator-only back-calculation on the realization of
    Equation \ref{eq:split}, and full-state conditioning.}
    \label{fig:YH-AW-18}
\end{figure}

\begin{figure}[H]
    \includegraphics[width=1\linewidth]{antiwindup_study/figures/YH-vs-H-cond-18.pdf}
    \caption{$H_\infty$ and Youla-H controllers, both with full-state conditioning:
    saturating step (left) and tracking error on a slow unsaturated ramp (right).}
    \label{fig:YH-vs-H-cond-18}
\end{figure}

% Add to each bibliography block:
% \bibitem{tustin1947method}
% Tustin, A. A Method of Analysing the Behaviour of Linear Systems in Terms of Time
% Series. \textit{Journal of the Institution of Electrical Engineers - Part IIA:
% Automatic Regulators and Servo Mechanisms} \textbf{1947}, \textit{94}, 130--142.
% \bibitem{hanus1987conditioning}
% Hanus, R.; Kinnaert, M.; Henrotte, J.-L. Conditioning Technique, a General
% Anti-windup and Bumpless Transfer Method. \textit{Automatica} \textbf{1987},
% \textit{23}, 729--739.
